% Source coverage: physical PDF pp. 375--376 (printed pp. 352--353), beginning at the Appendix heading and ending immediately before the Problems heading.
% Numbered equations: (21.A.1)--(21.A.11). Unnumbered displays, footnotes, figures, tables, citations, and centered dividers: none. Uncertainty: none.

\chapterappendix{}{General Unitarity Gauge}
\label{app:21.A}

In this appendix we shall show that in general spontaneously broken gauge
theories it is always possible to adopt a `unitarity' gauge in which the
Goldstone boson fields satisfy Eq.~\eqref{eq:21.4.30}:
\begin{equation}
 \sum_{rs}F_{rs}^2\xi_r e_{a s}=0.
 \tag{21.A.1}\label{eq:21.A.1}
\end{equation}

Using the exponential parameterization for all groups, we note first that
any element of $G$, in at least a finite neighborhood of the identity, may be
put in the form
\begin{equation}
 g=\exp\left(-i\sum_a\theta_a\mathcal T_a\right)
 \exp\left(i\sum_r\phi_r x_r\right)
 \exp\left(i\sum_i\mu_i t_i\right),
 \tag{21.A.2}\label{eq:21.A.2}
\end{equation}
with $\phi_r$ subject to the linear constraint that for all $a$
\begin{equation}
 \sum_{rs}F_{rs}^2\phi_r e_{a s}=0.
 \tag{21.A.3}\label{eq:21.A.3}
\end{equation}
This is easy to see when $g$ is infinitesimally close to the identity. Any such
$g$ may be written
\begin{equation}
 g=1+i\sum_r\phi_r^0x_r+i\sum_i\mu_i^0t_i,
 \tag{21.A.4}\label{eq:21.A.4}
\end{equation}
with $\phi_r^0$ and $\mu_i^0$ infinitesimal. Equivalently,
\begin{equation}
 g=1+i\sum_r\phi_rx_r+i\sum_i\mu_it_i
 -i\sum_a\theta_a\mathcal T_a,
 \tag{21.A.5}\label{eq:21.A.5}
\end{equation}
where $\theta_a$ is an arbitrary infinitesimal, and
\begin{equation}
 \phi_r(\theta)\equiv\phi_r^0+\sum_a\theta_a e_{a r},
 \tag{21.A.6}\label{eq:21.A.6}
\end{equation}
\begin{equation}
 \mu_i(\theta)\equiv\mu_i^0+\sum_a\theta_a e_{a i}.
 \tag{21.A.7}\label{eq:21.A.7}
\end{equation}
For any given $\phi_r^0$, we can choose $\theta_a$ to minimize the positive
quantity
\begin{equation}
 \sum_{rs}F_{rs}^2\phi_r(\theta)\phi_s(\theta).
 \tag{21.A.8}\label{eq:21.A.8}
\end{equation}
At this minimum, the quantity \eqref{eq:21.A.8} is stationary with respect to
variations of $\theta_a$, so $\phi_r(\theta)$ satisfies
Eq.~\eqref{eq:21.A.3}. For $\phi$, $\mu$, and $\theta$ infinitesimal,
Eq.~\eqref{eq:21.A.5} is the same as Eq.~\eqref{eq:21.A.2}, so we see that the
set of all $g$s of the form \eqref{eq:21.A.2} (with $\phi_r$ satisfying
Eq.~\eqref{eq:21.A.3}) includes all $g$s infinitesimally close to the identity.
It follows then from continuity that the same is true of all $g$s in at least
some finite neighborhood of the identity.

Next, consider the particular group element
\begin{equation}
 g=\gamma(\xi)=\exp\left(i\sum_r\xi_rx_r\right)
 \tag{21.A.9}\label{eq:21.A.9}
\end{equation}
and write it in the form \eqref{eq:21.A.2}:
\begin{equation}
 \gamma(\xi)
 =\exp\left(-i\sum_a\theta_a(\xi)\mathcal T_a\right)
 \gamma\bigl(\phi(\xi)\bigr)
 \exp\left(i\sum_i\mu_i(\xi)t_i\right),
 \tag{21.A.10}\label{eq:21.A.10}
\end{equation}
with $\phi_r(\xi)$ subject to Eq.~\eqref{eq:21.A.3}. This just says that the
gauge transformation\linebreak[4]
$\exp\left(i\sum_a\theta_a(\xi)\mathcal T_a\right)$
transforms $\xi_r$ into
\begin{equation}
 \xi'_r=\phi_r(\xi).
 \tag{21.A.11}\label{eq:21.A.11}
\end{equation}
Now dropping the prime, we have thus succeeded in constructing a gauge in
which $\xi_r$ satisfies Eq.~\eqref{eq:21.A.1}, as was to be shown.
